\documentclass[11pt]{article}

\usepackage[margin=1in]{geometry}
\usepackage[T1]{fontenc}
\usepackage[utf8]{inputenc}
\usepackage{lmodern}
\usepackage{amsmath,amssymb}
\usepackage{booktabs}
\usepackage{array}
\usepackage{enumitem}
\usepackage{longtable}
\usepackage{xcolor}
\usepackage{hyperref}
\usepackage{listings}

\hypersetup{
    colorlinks=true,
    linkcolor=blue,
    urlcolor=blue
}

\lstset{
    basicstyle=\ttfamily\small,
    breaklines=true,
    columns=fullflexible,
    keepspaces=true,
    frame=single
}

\title{Repository Map and Active Code Paths}
\author{}
\date{}

\begin{document}
\maketitle

\section{Current Entry Path}

The intended way to run the project is:
\begin{lstlisting}
python -m dev.main
\end{lstlisting}

The active call chain is:
\begin{lstlisting}
dev/main.py
-> dev/experiments/generalized_study.py
-> dev/experiments/study/
\end{lstlisting}

\texttt{generalized\_study.py} is only a thin compatibility wrapper. The real orchestration lives in the \texttt{study} package.

\section{Where the Important Logic Lives}

\subsection{\texttt{dev/master\_config.py}}

This is the single user-editable configuration file. It defines:
\begin{itemize}[leftmargin=2em]
    \item which branch is active via \texttt{MAP\_TYPE},
    \item whether raw MAPF data is recomputed via \texttt{recompute\_MAPF},
    \item which presentation artifact family is regenerated via \texttt{to\_generate},
    \item the global limit on consecutive failed paired sampling attempts,
    \item the global \texttt{enhanced\_CBS} toggle,
    \item the global \texttt{compact\_clustering} toggle,
    \item each branch's editable \texttt{ECBS\_suboptimality} value,
    \item each branch's low-level guidance toggles \texttt{true\_static\_shortest\_path\_distance} and \texttt{tight\_time\_horizon},
    \item campus image paths and semantic-color settings,
    \item all other branch-specific numeric parameters and toggles.
\end{itemize}

\subsection{\texttt{dev/experiments/branch\_specs.py}}

This file converts user config dictionaries into normalized branch specifications. In practice, many parts of the runner depend on the \texttt{BranchSpec} dataclass rather than directly reading raw config dictionaries. The branch spec now also records the effective solver family metadata, the active clustering style, campus lane-variant metadata when relevant, and the two branch-level low-level guidance toggles, so the output directory keeps track of which heuristic and horizon policy were active for the run.

\subsection{\texttt{dev/experiments/study/}}

This is the active experiment engine. Its submodules are separated by responsibility:
\begin{itemize}[leftmargin=2em]
    \item \texttt{models.py}: dataclasses for run configurations, run records, aggregates, and dynamic state,
    \item \texttt{preparation.py}: branch-aware setup logic,
    \item \texttt{runtime.py}: solver execution, solver-mode dispatch, and result-category conversion,
    \item \texttt{aggregation.py}: condition-level summaries,
    \item \texttt{io\_utils.py}: JSON/CSV writing and immediate console/log flushing,
    \item \texttt{logging\_utils.py}: console/log formatting,
    \item \texttt{metrics\_data\_store.py}: Results-ready configuration packages, descriptive summaries, paired comparisons, reader guides, and the project-level data dictionary,
    \item \texttt{raw\_data\_store.py}: branch-local persisted raw MAPF storage under \texttt{dev/raw\_mapf\_data/}, organized as a manifest plus metadata and per-condition files rather than one monolithic pickle,
    \item \texttt{visualization.py}: selection and rendering of final Pillow exports, with selection driven by the current visualization settings even when the frame ingredients come from persisted raw MAPF data,
    \item \texttt{orchestrator.py}: the main control flow.
\end{itemize}

\subsection{\texttt{dev/mapf/}}

This folder contains the actual planning code:
\begin{itemize}[leftmargin=2em]
    \item static CBS/ECBS and low-level A*,
    \item time-expanded CBS/ECBS and low-level A*,
    \item low-level guidance helpers in \texttt{low\_level\_guidance.py} for exact static-distance heuristics and tighter horizon support,
    \item conflict metrics,
    \item frame builders and image loggers.
\end{itemize}

\subsection{\texttt{dev/maps/} and \texttt{dev/navigation/}}

These modules implement the map representation, mapping overlays, cyclic cleanup, connectivity restoration, and neighbor-generation semantics over composite grids.

\subsection{\texttt{dev/inputs/}}

Image-based branches depend on this folder. The current dynamic input logic supports both ordinary thresholded port loading and campus semantic-color loading. The assets currently referenced by \texttt{master\_config.py} are \texttt{dynamic\_port/port\_map.png}, \texttt{campus\_area\_1.png}, \texttt{campus\_area\_2.png}, and \texttt{campus\_area\_3.png}. Static and dynamic variants of the same port or campus area share the same source image.

\section{Legacy vs.\ Active Code}

The project still contains older branch-specific pipeline modules such as:
\begin{itemize}[leftmargin=2em]
    \item \texttt{dev/experiments/static\_artificial/pipeline.py},
    \item \texttt{dev/experiments/dynamic\_port/pipeline.py}.
\end{itemize}

These are useful references for branch intent, but the current top-level execution path does \emph{not} go through them. When updating the project, changes that affect current behavior should usually be made under:
\[
\texttt{master\_config.py},\quad
\texttt{experiments/branch\_specs.py},\quad
\texttt{experiments/study/},\quad
\texttt{mapf/}.
\]

\section{Output Contract}

Main-experiment outputs are written under \texttt{dev/outputs\_main/}. The active output families are:
\begin{itemize}[leftmargin=2em]
    \item \texttt{metrics\_data\_inspection/}: author-facing evaluation text for each exact configuration;
    \item \texttt{metrics\_data/}: independent Results-ready per-configuration packages plus the project-level data dictionary;
    \item \texttt{frame\_by\_frame/}: designated successful trajectory packages;
    \item \texttt{terminal\_logs/}: one \texttt{raw\_data.log} or \texttt{visualization.log} per exact configuration;
    \item \texttt{visualization/}: Pillow frames generated from the saved trajectory packages.
\end{itemize}

The exact map hierarchy is retained beneath the metrics, frame-by-frame, terminal-log, and visualization trees. The main experiment no longer has a graph-generation mode or a PNG metric-output folder. Reference-comparison outputs remain separate under \texttt{dev/outputs\_ref\_comparison/}.

\end{document}
